\reading{LTR-7}{Monitoring Liquidity}{STRIKE FIRST}{2}

\begin{lrkeybox}[Learning objectives — official 2026 spine wording]
\begin{enumerate}[label=\textbf{\alph*.}]
\item Differentiate between deterministic and stochastic cash flows and provide examples of each.
\item Describe and identify examples of liquidity options and explain the impact of liquidity options on a bank's liquidity position and its liquidity management process.
\item Describe and apply the concepts of liquidity risk, funding cost risk, liquidity generation capacity, expected liquidity, and cash flow at risk.
\item Interpret the term structure of expected cash flows and cumulative cash flows.
\item Discuss the impact of available asset transactions on cash flows and liquidity generation capacity.
\end{enumerate}
{\footnotesize\color{FRMGrey}Source: Antonio Castagna and Francesco Fede, \textit{Measuring and Managing Liquidity Risk}, Chapter 6.}
\end{lrkeybox}

\begin{lrtrapbox}[The label collision, and the ordering]
The Lecture layer labels this reading \texttt{LO 65.b--e} --- the same \texttt{65} prefix the Crash layer uses for \textbf{LTR-2 Liquidity and Leverage}. It also runs them out of order (b, d, c, e), and it has no coverage of objective \textbf{a}, which exists only in the Crash layer. Everything below is renumbered and resequenced to the spine.
\end{lrtrapbox}

% =====================================================================
\lo{LTR-7 a}{Differentiate between deterministic and stochastic cash flows and provide examples of each.}

\begin{lrdefbox}[Two dimensions, two labels]
Cash flows at a point in time can be identified on \textbf{two dimensions: time and amount}.
\begin{itemize}
\item In terms of \textbf{time}, cash flows are \term{deterministic} if they occur at known times, and \term{stochastic} if they occur randomly.
\item In terms of \textbf{amount}, cash flows whose amounts are known are \term{deterministic}, and cash flows with unknown amounts are \term{stochastic}.
\end{itemize}
\end{lrdefbox}

\begin{lrtrapbox}[A cash flow can be deterministic on one dimension and stochastic on the other]
The classification is not a single binary. A fixed-coupon bond pays a known amount at a known date --- deterministic on both. A floating-rate coupon pays an unknown amount at a known date --- deterministic in time, stochastic in amount. A deposit withdrawal is stochastic on both. An option to classify each item on \emph{one} dimension only, or a stem that supplies only the timing and asks for the full classification, is testing exactly this.
\end{lrtrapbox}

% =====================================================================
\lo{LTR-7 b}{Describe and identify examples of liquidity options and explain the impact of liquidity options on a bank's liquidity position and its liquidity management process.}

\begin{lrdefbox}[Liquidity option]
The right of a holder to receive cash from, or to give cash to, the bank at predefined times and terms.
\end{lrdefbox}

\begin{lrkeybox}[How a liquidity option differs from a standard option]
\begin{itemize}
\item \textbf{Standard options} traded in financial markets are exercised \textbf{when there is a profit}.
\item A \textbf{liquidity option} is exercised \textbf{because of the cash flows produced after exercise, even if it makes a loss}.
\end{itemize}
This is the whole discriminator. A liquidity option is not a rational-exercise instrument, which is why it cannot be modelled with standard option pricing and why behavioural models are needed instead (see LTR-7 d).
\end{lrkeybox}

\subsection{Examples}

\begin{itemize}
\item Lines of credit that allow customers to withdraw funds at their discretion.
\item The ability to withdraw funds from savings accounts with minimal notice.
\end{itemize}

Customers may use these options for various reasons, such as taking advantage of favourable interest rates or simply for convenience.

\subsection{Two ways exercise affects the bank}

\begin{enumerate}
\item \textbf{Balance sheet.} The bank's balance sheet is affected by the amount repaid or withdrawn. Withdrawals reduce the bank's \emph{cash} and the corresponding \emph{deposit liability}; repayments increase cash.
\item \textbf{Financial impact.} Positive or negative, depending on the spread between the initial contract's interest rates, spread and market values, compared with conditions at the time of exercise.
\end{enumerate}

To manage these impacts, banks may use cash reserves, liquid assets, or access to credit lines as hedging tools.

\begin{lrtrapbox}[Whose profit? State the perspective before you sign the impact]
The two worked directions in the source notes are stated from the \emph{customer's} point of view, and they invert from the bank's:
\begin{itemize}
\item A customer draws a credit line at a \textbf{fixed rate below the current market rate}. Good for the customer; the bank must fund the drawdown at market and earns less than market, so the impact on the \textbf{bank is negative}.
\item A customer withdraws from a \textbf{high-interest savings account when market rates have fallen below the contracted rate}. Good for the customer to leave; the bank was paying above market, so the withdrawal is \textbf{positive for the bank}.
\end{itemize}
The objective asks about the impact \emph{on the bank}. Name the actor, then the direction. See the corrections record.
\end{lrtrapbox}

% =====================================================================
\lo{LTR-7 c}{Describe and apply the concepts of liquidity risk, funding cost risk, liquidity generation capacity, expected liquidity, and cash flow at risk.}

\subsection{The four risk definitions}

\begin{center}
\small
\begin{tabularx}{\textwidth}{@{}p{3.7cm} X@{}}
\toprule
\rowcolor{FRMSoft}\textbf{Concept} & \textbf{Definition} \\
\midrule
\textbf{Quantitative liquidity risk} & A bank's challenge of receiving \emph{lower-than-expected cash flows} in the future, making it difficult to meet upcoming payment obligations. \\
\addlinespace
\textbf{Market liquidity risk} & A \textbf{subset} of quantitative liquidity risk. Occurs when assets cannot be swiftly sold at a fair price, leading to lower-than-expected cash inflows from sales. \\
\addlinespace
\textbf{Funding cost risk} & Also known as the \emph{cost of liquidity risk}. Arises when a bank has to bear \emph{higher-than-expected costs}, such as the spread over the risk-free rate, to obtain funds from available liquidity sources. \\
\addlinespace
\textbf{Liquidity risk} (the umbrella) & Encompasses \textbf{both} the quantitative and the cost dimensions. Represents economic losses resulting from variations in cash flows, and in available cash on a specific date, differing from expectations. \\
\bottomrule
\end{tabularx}
\end{center}

\begin{lrtrapbox}[Quantity versus cost, and the nesting]
\textbf{Quantitative} risk is about \emph{how much cash arrives}; \textbf{funding cost} risk is about \emph{what the cash costs}. Market liquidity risk sits \emph{inside} quantitative risk, not alongside it, and overall liquidity risk sits above both. A statement placing market liquidity risk parallel to funding cost risk has broken the hierarchy --- a classic scope corruption.
\end{lrtrapbox}

\subsection{Liquidity generation capacity}

\begin{lrdefbox}[Liquidity generation capacity (LGC)]
The main tool a bank can use to handle the \textbf{negative entries of the TSECCF}. The ability of a bank to generate positive cash flows, \textbf{beyond contractual ones}, from the sources of liquidity available on the balance sheet and off the balance sheet at a given date.
\end{lrdefbox}

\begin{center}
\small
\begin{tabularx}{\textwidth}{@{}p{4.3cm} X@{}}
\toprule
\rowcolor{FRMSoft}\textbf{How LGC manifests} & \textbf{What it means} \\
\midrule
\textbf{1. Balance sheet expansion} \newline with secured or unsecured funding &
Borrowing through an increase of deposits, typically in the interbank market; withdrawal of credit lines; issuance of new bonds. \\
\addlinespace
\textbf{2. Balance sheet shrinkage} &
Operated by selling assets, starting from the more liquid ones such as Treasury bonds, corporate bonds and stocks. \\
\addlinespace
\textbf{Repo transactions} &
Can be considered separately from the other two cases and labelled \textbf{``balance sheet neutral''}. \\
\bottomrule
\end{tabularx}
\end{center}

\begin{lrtrapbox}[Repo is the third category, not a member of the first two]
Expansion grows the balance sheet, shrinkage contracts it, and repo does neither --- it is explicitly \textbf{balance-sheet neutral}. An option filing repo under ``expansion with secured funding'' is the intended sibling swap, and it reads well because a repo \emph{is} secured funding.
\end{lrtrapbox}

\subsection{Cash flow at risk}

The fact that cash flows are stochastic suggests the need for some measure related to their volatility.

\begin{lrfmlbox}[Cash flow at risk, both signs]
\[ \text{cfaR}^{+} \;=\; \text{Expected inflow} - \text{Actual inflow} \qquad\qquad
   \text{cfaR}^{-} \;=\; \text{Actual outflow} - \text{Expected outflow} \]
\vk{Positive cfaR is defined at confidence $\alpha$ and measures the shortfall when you \emph{expect more and get less}. Negative cfaR is defined at $1-\alpha$ and measures the excess when you \emph{pay more than expected}. Both are stated as positive numbers representing an adverse surprise.}
\end{lrfmlbox}

\begin{lrexbox}[The two worked cases the notes give]
\textbf{Positive cash flow at risk.} Expected inflow \$100m, actual inflow \$90m:
\[ \text{cfaR}^{+} = 100 - 90 = \mathbf{\$10\text{ million}} \qquad \text{(maximum shortfall at 95\%)} \]
\textbf{Negative cash flow at risk.} Expected outflow \$12m, actual outflow \$15m:
\[ \text{cfaR}^{-} = 15 - 12 = \mathbf{\$3\text{ million}} \qquad \text{(minimum excess at 5\%)} \]
\tcblower
\textbf{Read the subtraction order.} The two cases subtract in \emph{opposite} orders, and that is not sloppiness --- it is because an adverse surprise on an inflow is a shortfall (expected minus actual) while an adverse surprise on an outflow is an excess (actual minus expected). Both produce a positive number. The source notes write the first one as $90 - 100$, which gives $-10$; corrected here. See the corrections record.
\end{lrexbox}

\begin{lrtrapbox}[The sign convention is the trap]
Both quantities are reported as \textbf{positive} magnitudes of an adverse move. If a calculation returns a negative cfaR, the subtraction order is wrong. This is a \S5B.3 sign trap, and $-10$ will be an offered option.
\end{lrtrapbox}

\subsection{The term structures built from cfaR}

\begin{lrfmlbox}[Term structures of unexpected cash flows and liquidity at risk]
\[ \text{TSCF}^{+}_{\alpha}(t_0,t_b) = \bigl\{ \text{cfaR}^{+}_{\alpha}(t_0,t_0),\; \text{cfaR}^{+}_{\alpha}(t_0,t_1),\;\ldots,\; \text{cfaR}^{+}_{\alpha}(t_0,t_b) \bigr\} \]
\[ \text{TSCF}^{-}_{1-\alpha}(t_0,t_b) = \bigl\{ \text{cfaR}^{-}_{1-\alpha}(t_0,t_0),\; \ldots,\; \text{cfaR}^{-}_{1-\alpha}(t_0,t_b) \bigr\} \]
\[ \text{TSLaR}_{1-\alpha}(t_0,t_b) = \bigl\{ \text{cf}_{1-\alpha}(t_0,t_i) - \text{TSECCF}(t_0,t_i) - \text{TSCLGC}(t_0,t_i) \bigr\}_{i=1,\ldots,b} \]
\vk{$\alpha$ = confidence level; $t_0$ = reference date; $t_b$ = terminal date. The \term{term structure of liquidity at risk (TSLaR)} is the collection of unexpected cash flows at each date over a period --- the difference between the average and minimum levels of cash flows --- net of both the cumulated expected cash flows and the cumulated liquidity generation capacity.}
\end{lrfmlbox}

\begin{lrkeybox}[The source notes' own guidance on these formulas]
The source notes state plainly: \textbf{``Don't memorise the formulas. Not needed for the exam.''} That is the notes' judgement and it is retained here. Note what the spine actually asks for, though: it says \emph{describe and apply the concepts}. The concepts --- what LGC is, what cfaR measures, why TSLaR subtracts both TSECCF and TSCLGC --- are examinable even if the subscript notation is not.
\end{lrkeybox}

% =====================================================================
\lo{LTR-7 d}{Interpret the term structure of expected cash flows and cumulative cash flows.}

\subsection{Two classes of factor produce cash flows}

\begin{center}
\small
\begin{tabularx}{\textwidth}{@{}p{3.4cm} X p{2.4cm}@{}}
\toprule
\rowcolor{FRMSoft}\textbf{Class} & \textbf{Definition} & \textbf{Term structures} \\
\midrule
\textbf{Causes of liquidity} &
All factors referring to existing and forecast future contracts originated by the \textbf{ordinary business activity} of a financial institution. Cash flows generated by the causes of liquidity can be \textbf{both positive and negative}. &
TSECF, TSECCF \\
\addlinespace
\textbf{Sources of liquidity} &
All factors capable of generating \textbf{positive} cash flows to manage and hedge liquidity risk, and which can be disposed of promptly by the bank. These determine the \textbf{liquidity generation capacity (LGC)}. &
TSLGC, TSCLGC \\
\bottomrule
\end{tabularx}
\end{center}
\figcap{The asymmetry is the point: \textbf{causes} run in both directions, \textbf{sources} only generate positive flows. That is why LGC is the instrument for fixing a negative TSECCF and never the thing that creates one.}

\begin{lrdefbox}[The four term structures]
\begin{itemize}
\item \term{TSECF} --- the \textbf{term structure of expected cash flows}, also known as the \textbf{maturity ladder}: the collection by date of positive and negative expected cash flows, up until the expiration of the contract with the longest maturity.
\item \term{TSECCF} --- the \textbf{term structure of expected cumulated cash flows}: the collection of cumulative expected cash flows in date order.
\item \term{TSLGC} --- the term structure of LGC: the liquidity at a given time that can be generated by sources of liquidity up to a terminal time $t_b$. All values are expected values.
\item \term{TSCLGC} --- the term structure of cumulated LGC: the cumulative liquidity that can be generated by sources of liquidity up to $t_b$.
\end{itemize}
\end{lrdefbox}

\begin{lrfmlbox}[TSECF and cumulated LGC in the source's notation]
\[ \text{TSECF}(t_0,t_b) = \bigl\{ \text{cf}^{+}_{e}(t_0,t_0),\, \text{cf}^{-}_{e}(t_0,t_0),\, \text{cf}^{+}_{e}(t_0,t_1),\, \text{cf}^{-}_{e}(t_0,t_1),\, \ldots,\, \text{cf}^{+}_{e}(t_0,t_b),\, \text{cf}^{-}_{e}(t_0,t_b) \bigr\} \]
\[ \text{TSCLGC}(t_0,t_b) = \Bigl\{ \textstyle\sum_{i=0}^{1}\text{TSLGC}(t_0,t_i),\; \sum_{i=0}^{2}\text{TSLGC}(t_0,t_i),\; \ldots,\; \sum_{i=0}^{b}\text{TSLGC}(t_0,t_i) \Bigr\} \]
\vk{$\text{cf}^{+}_{e}$ and $\text{cf}^{-}_{e}$ = expected positive and negative cash flows; $t_0$ = reference date; $t_i$ = the $i$-th future date; $t_b$ = terminal date. Note that TSECF holds the positive and negative flows \emph{separately} at each date, rather than netting them.}
\end{lrfmlbox}

\subsection{How the term structures are built}

\begin{itemize}
\item To calculate TSECF and TSECCF, banks start with a \textbf{balance sheet} showing assets, liabilities and equities.
\item Assuming no liquidity options within deposits and no default risk in assets, \textbf{expiring assets generate positive cash flows} while \textbf{expiring liabilities result in negative cash flows}.
\item Assets and liabilities should be organized by their expiration dates, and interest payments from these assets and liabilities need to be considered.
\end{itemize}

\begin{lrexbox}[TSECCF for King Bank]
At the start of Year 1, King Bank's balance sheet shows assets of \$125 million, liabilities of \$75 million and equity of \$50 million. At the end of Year 1, assets worth \$25 million expire. At the end of Year 2, liabilities worth \$15 million expire. Assets earn 6.00\% and liabilities bear 5.00\%. Calculate TSECCF for Years 1 and 2.
\tcblower
\textbf{Year 1.} Three flows: the expiring asset, the interest earned on the full asset base, and the interest paid on the full liability base.
\begin{align*}
\text{TSECCF}(1) &= \underbrace{25}_{\text{expiring asset}} \;+\; \underbrace{0.06 \times 125 = 7.50}_{\text{interest earned}} \;-\; \underbrace{0.05 \times 75 = 3.75}_{\text{interest paid}} \\
&= \mathbf{\$28.75\text{ million}}
\end{align*}

\textbf{Year 2.} Start from the Year 1 cumulative figure. Assets are now \$100 million, because \$25 million expired; liabilities are still \$75 million until they expire at the \emph{end} of Year 2, so the full 5\% is still paid.
\begin{align*}
\text{TSECCF}(2) &= \underbrace{28.75}_{\text{from Year 1}} \;+\; \underbrace{0.06 \times 100 = 6.00}_{\text{interest earned}} \;-\; \underbrace{15}_{\text{expiring liability}} \;-\; \underbrace{0.05 \times 75 = 3.75}_{\text{interest paid}} \\
&= \mathbf{\$16.00\text{ million}}
\end{align*}

\textbf{What this means.} TSECCF is \emph{cumulative}, so Year 2 builds on Year 1 rather than standing alone. The figure stays positive across both years, which is what the Treasury wants: positive cash flows covering negative ones.
\end{lrexbox}

\begin{lrtrapbox}[Three ways this example is corrupted]
\begin{itemize}
\item \textbf{Interest base in Year 2.} Assets earn on \$100m (the \$25m expired), but liabilities still pay on the full \$75m, because they expire at the \emph{end} of Year 2. Using \$60m for the liability base gives a clean wrong answer.
\item \textbf{Forgetting the carry-forward.} Year 2 \emph{non-cumulative} flow is $6.00 - 15 - 3.75 = -12.75$. That is the TSECF figure, not the TSECCF figure, and it is the classic intermediate-result distractor.
\item \textbf{Sign on the expiring item.} An expiring \textbf{asset} is a positive flow; an expiring \textbf{liability} is negative. Treating the expiring liability as a positive flow turns \$16.00m into \$46.00m.
\end{itemize}
\end{lrtrapbox}

\subsection{Key elements of the model}

\begin{enumerate}
\item Cash flows are often \textbf{stochastic} due to their linkage to market indices.
\item \textbf{Credit models} need to consider default risk and the correlations that inevitably exist between bank counterparties.
\item \textbf{Behavioural models} can be used to adjust cash flows for liquidity options.
\item \textbf{New business cash flows} should also be considered, as well as rollovers of maturing liabilities and new bond issuances to fund asset increases.
\end{enumerate}

\begin{lrkeybox}[What the Treasury Department is monitoring, and why]
\begin{itemize}
\item The goal is to monitor TSECF and TSECCF, with the \textbf{latter ideally being positive at all times}, implying that positive cash flows can cover negative cash flows.
\item Because TSECF forecasts \emph{expected} values of cash flows, TSECCF only captures expected values as well.
\item Because of the timing of asset and liability maturities, \textbf{there may be periods of negative cumulated cash flows}.
\item \textbf{A negative balance in TSECCF could be a precursor to insolvency and bankruptcy}, which is why maintaining a positive balance matters.
\end{itemize}
\end{lrkeybox}

\subsection{Building the TSLGC is harder than building the TSECF}

\begin{itemize}
\item The assumptions made for a given period affect other periods.
\item The TSLGC is intertwined with the TSECF, and therefore with the TSECCF.
\item Most of the problems in building the TSLGC are caused by \textbf{unencumbered assets}, most of which are AFS bonds. The bank must keep track of how many are either sold or repoed out, and the liquidity amount extractable from these assets depends on several risk factors.
\end{itemize}

% =====================================================================
\lo{LTR-7 e}{Discuss the impact of available asset transactions on cash flows and liquidity generation capacity.}

\begin{lrdefbox}[The term structure of available assets (TSAA)]
The TSAA shows how many single securities, or all of them, are available for inclusion in LGC. According to the type of operation --- selling or repo --- the \textbf{price} and/or the \textbf{haircut} are factors that need to be considered.
\end{lrdefbox}

The treasury department's transactions fall into three families: repo and reverse repo; buy/sellback and sell/buyback; and lend and borrow.

\subsection{The impact grid}

\begin{center}
\scriptsize
\begin{longtable}{@{}p{3.5cm} p{2.3cm} p{2.9cm} p{5.2cm}@{}}
\toprule
\rowcolor{FRMSoft}\textbf{Transaction} & \textbf{TSAA} & \textbf{TSECF} & \textbf{LGC} \\
\midrule
\endfirsthead
\toprule
\rowcolor{FRMSoft}\textbf{Transaction} & \textbf{TSAA} & \textbf{TSECF} & \textbf{LGC} \\
\midrule
\endhead
\textbf{Asset purchase} & Up & Down & No change \\
\textbf{Asset sale} & Down & Up & Yes \\
\addlinespace
\textbf{Buy/sellback} --- at the start & Up & Down & No change \\
\textbf{Buy/sellback} --- at the end & Down & Up & No (exception) \\
\addlinespace
\textbf{Sell/buyback} --- at the start & Down & No change & Yes \\
\textbf{Sell/buyback} --- at the end & Up & Down by interest & No change \\
\addlinespace
\textbf{Securities lending} & Down & Up by fees & No change \\
\textbf{Securities borrowing} & Up & Down by fees & No change \\
\addlinespace
\textbf{Repo} --- asset leg (owned, transferred) & Down (NP) & No change & Yes (repo $-$ haircut) \\
\textbf{Repo} --- counterparty side & Up (NP) & Down by interest & Zero (repo $-$ haircut) \\
\addlinespace
\textbf{Reverse repo} --- own side & Up (NP) & Down (repo $-$ haircut) & No change \\
\textbf{Reverse repo} --- counterparty side & Down (NP) & Up by cash $+$ interest & No change \\
\bottomrule
\end{longtable}
\end{center}
\figcap{``NP'' marks a transfer where the asset is \emph{owned but transferred} --- no purchase. The grid's logic: TSAA tracks what the bank \emph{holds and can pledge}; TSECF tracks \emph{cash moving}; LGC records only whether the transaction \emph{generated} liquidity. Note the asymmetry that catches people: a sell/buyback generates LGC at the start and nothing at the end, while a buy/sellback generates nothing at either point.}

\subsection{The three transaction structures}

\begin{center}
\begin{tikzpicture}[
  bx/.style={draw=FRMNavy, fill=PaleNavy, rounded corners=1pt, align=center,
             font=\scriptsize\sffamily\bfseries, minimum width=2.1cm, minimum height=0.7cm},
  ar/.style={-{Stealth[length=4.5pt]}, thick},
  lb/.style={font=\tiny\sffamily, fill=white, inner sep=1.5pt}
]
% repo
\node[font=\scriptsize\sffamily\bfseries, anchor=west] at (0.15,2.35) {Repo (bank is the cash borrower)};
\node[bx] (b1) at (1.3,0.9) {Bank};
\node[bx] (c1) at (5.5,0.9) {Counterparty};
\draw[ar, FRMRed] (b1.north east) to[bend left=26] node[lb,pos=0.5]{Asset} (c1.north west);
\draw[ar, FRMGreen] (c1.south west) to[bend left=26] node[lb,pos=0.5]{Cash} (b1.south east);
% reverse repo
\node[font=\scriptsize\sffamily\bfseries, anchor=west] at (8.15,2.35) {Reverse repo (bank is the cash lender)};
\node[bx] (b2) at (9.3,0.9) {Bank};
\node[bx] (c2) at (13.5,0.9) {Counterparty};
\draw[ar, FRMRed] (b2.north east) to[bend left=26] node[lb,pos=0.5]{Cash $+$ int.} (c2.north west);
\draw[ar, FRMGreen] (c2.south west) to[bend left=26] node[lb,pos=0.5]{Asset} (b2.south east);
\node at (0,-0.35) {};
\end{tikzpicture}
\figcap{Which way the \emph{asset} moves is the only thing that has to be remembered; the cash always moves the other way. In a repo the bank gives up the asset and takes in cash, so LGC rises by the repo value net of the haircut. In a reverse repo the bank hands over cash, so no liquidity is generated.}
\end{center}

\begin{lrexbox}[Buy/sellback --- the bank buys now and sells back later]
The bank buys 40,000 bonds and sells them back at the forward price. At the start of the contract the bank pays \textbf{500,000}. At the end of the contract the bank sells the bond back at the contract price; the sum it receives also includes accrued interest, \textbf{501,000}.
\tcblower
\textbf{At the start.}
\begin{itemize}
\item TSECF $= (-500{,}000)$ --- cash goes out.
\item Bonds available increase to 40,000 in the TSAA.
\item The TSCLGC is \textbf{not modified}.
\end{itemize}
\textbf{At the end.}
\begin{itemize}
\item TSECF $= (+501{,}000)$ --- cash comes back, including accrued interest.
\item The TSAA shows a reduction of the available quantity back to zero bonds.
\end{itemize}
\textbf{Net cash effect:} $501{,}000 - 500{,}000 = \mathbf{+1{,}000}$, which is the accrued interest earned over the life of the contract. The bank bought and returned the same bond; the only economic residue is the interest.
\end{lrexbox}

\begin{lrexbox}[Sell/buyback --- the bank sells now and buys back later]
At the start the bank sells the bond and receives \textbf{307,200}. At the end the bank buys the bond back and pays \textbf{307,350}.
\tcblower
\textbf{At the start.}
\begin{itemize}
\item TSCLGC $= (+307{,}200)$ --- this transaction \emph{generates} liquidity.
\item The TSAA indicates a reduction of the available quantity.
\end{itemize}
\textbf{At the end.}
\[ 307{,}350 - 307{,}200 = \mathbf{150} \]
\begin{itemize}
\item The difference of \textbf{150} enters the TSECF on the end date of the sell/buyback.
\item The TSAA increases back to its starting quantity of 70,000 bonds.
\end{itemize}
\textbf{The contrast worth holding.} In the buy/sellback, the full cash amounts hit the TSECF at both ends and LGC never moves. In the sell/buyback, the initial proceeds hit \emph{TSCLGC} and only the \textbf{net difference} hits the TSECF at the end. Same two legs, opposite direction, and the accounting routes differ.
\end{lrexbox}

\begin{lrkeybox}[Securities lending versus securities borrowing]
\begin{itemize}
\item \textbf{Securities lending} is similar to a sell/buyback, except that \textbf{only possession passes} to the counterparty --- no cash is paid by the counterparty to the bank. The bank receives fee payments for the asset, along with interest from the counterparty at contract expiration, which positively impacts TSECF and TSECCF.
\item \textbf{Securities borrowing} is similar to a buy/sellback, except that possession passes \emph{to} the bank and no cash is paid by the bank. The counterparty receives payments for the asset. The only impact on TSECF and TSECCF is the interest payment made by the bank at contract expiration.
\end{itemize}
In both cases \textbf{LGC does not change}, because no cash was generated --- which is precisely what separates them from the repo pair.
\end{lrkeybox}

\begin{lrtrapbox}[Consolidated trap summary — LTR-7]
\begin{itemize}
\item \textbf{Scope.} Deterministic/stochastic is a classification on \emph{two} dimensions, time and amount. An instrument can be one on each.
\item \textbf{Definition.} A standard option is exercised \emph{for profit}; a liquidity option is exercised \emph{for the cash flow, even at a loss}.
\item \textbf{Role.} The financial impact of an exercised liquidity option must be signed \emph{from the bank's side}, and it is the opposite of the customer's.
\item \textbf{Scope.} Market liquidity risk is a \emph{subset} of quantitative liquidity risk, not a parallel category; funding cost risk is the cost dimension; liquidity risk is the umbrella over both.
\item \textbf{Sibling.} LGC via balance sheet \emph{expansion} vs \emph{shrinkage}, with repo as a separate \textbf{balance-sheet-neutral} third case.
\item \textbf{Sign.} $\text{cfaR}^{+} = $ expected $-$ actual; $\text{cfaR}^{-} = $ actual $-$ expected. Both reported positive.
\item \textbf{Intermediate result.} In a TSECCF question the single-period TSECF figure ($-12.75$m in the worked example) is the offered distractor against the cumulative \$16.00m.
\item \textbf{Calculation.} In a multi-period TSECCF, the interest base changes only \emph{after} the item expires. Assets fall to \$100m in Year 2; liabilities stay at \$75m.
\item \textbf{Polarity.} Expiring \textbf{assets} give positive cash flows; expiring \textbf{liabilities} give negative ones.
\item \textbf{Sibling.} Sell/buyback generates LGC at the start; buy/sellback never does. Securities lending and borrowing never do either.
\end{itemize}
\end{lrtrapbox}
